\section{Implementation}
\label{sec:implementation}

% Target: ~2 two-column pages (~4 single-column equivalent)
% Condense report sections/implementation.tex. Focus on concrete mechanisms,
% data structures, and the domain-switch path. Use code snippets sparingly
% (IEEE page budget is tight).

\subsection{API Extensions}
\label{sub:api}

% [~0.3 column]
%   - FreeRTOS-MPU creates tasks via `xTaskCreateRestricted()`. This
%     function receives a `TaskParameters_t` struct containing stack,
%     priority, and MPU region descriptors.
%   - Extension: add `DomainParameters_t` with fields `uxStartSlice` and
%     `uxNumSlices`. If `uxNumSlices == 0`, the task is created in the
%     caller's domain (no domain carving).
%   - If a new domain is created, the slices are carved from Domain 0.
%     Domain 0 is the implicit initial domain.
%   - Short code snippet showing the extended struct definition (4--5 lines).

\subsection{Configuration Flags}
\label{sub:config}

% [~0.2 column]
%   - Three compile-time flags in `FreeRTOSConfig.h`:
%     \begin{itemize}
%       \item \texttt{configENABLE\_DOMAINS}: global enable/disable (0 or 1).
%       \item \texttt{configNUM\_TIME\_SLICES}: total number of slices in the
%         scheduling timeline.
%       \item \texttt{configNUM\_TICKS\_PER\_SLICE}: OS ticks per slice.
%     \end{itemize}
%   - When domains are disabled, the kernel compiles to standard FreeRTOS
%     with zero overhead.

\subsection{Domain Switching}
\label{sub:domain_switching}

% [~1 column]
%   - Hook into `xTaskIncrementTick()` (called from machine timer interrupt).
%     Each tick increments the effective tick counter. Domain tick =
%       effective\_tick / configNUM\_TICKS\_PER\_SLICE (integer division,
%       wrapping at configNUM\_TIME\_SLICES).
%   - Maintain a `uxCurrentDomainIndex` tracking the active domain for the
%     current slice. On a domain-tick change, look up `xDomains[domain\_tick]`
%     to determine the target domain.
%   - If `target\_domain != current\_domain`:
%     \begin{enumerate}
%       \item Invoke `temporal\_fence\_t()` to flush on-core
%         microarchitectural state.
%       \item Save the current task's context (using FreeRTOS's existing
%         `portSAVE\_CONTEXT`).
%       \item Select the highest-priority ready task from the target domain's
%         ready list.
%       \item Restore that task's context.
%     \end{enumerate}
%   - `temporal\_fence\_t()` is implemented as inline assembly:
%     \begin{verbatim}
%       asm volatile("fence.t" ::: "memory");
%     \end{verbatim}
%   - On QEMU, this is a no-op (no microarchitectural state to flush) but
%     the instruction is recognised and the execution path is exercised.

\subsection{Per-Domain Scheduling State}
\label{sub:per_domain_state}

% [~0.5 column]
%   - FreeRTOS tracks scheduling state in global variables:
%     `pxCurrentTCB`, task ready lists (`pxReadyTasksLists`),
%     `uxCurrentNumberOfTasks`, etc.
%   - Domain extension: convert these to arrays indexed by domain ID.
%     `pxCurrentTCBs[domainID]` replaces `pxCurrentTCB`;
%     `pxReadyTasksLists[domainID][priority]` replaces
%     `pxReadyTasksLists[priority]`.
%   - The tick-interrupt handler selects `pxCurrentTCBs[target\_domain]`
%     after the domain switch.
%   - Within-domain scheduling is unchanged: `vTaskSwitchContext()` is
%     parameterised by domain and operates on that domain's ready lists.

\subsection{Cross-Domain Queues}
\label{sub:queues}

% [~0.3 column]
%   - Each TCB stores `uxDomainID`. When a queue `xQueueSend()` or
%     `xQueueReceive()` unblocks a task, the task's TCB domain ID is used
%     to place it in the correct domain's ready list.
%   - No changes to the queue data structure itself; only the insertion
%     logic in `xTaskRemoveFromEventList()` is modified to use the task's
%     `uxDomainID`.
%   - Acknowledged limitation: tasks that share a queue can still infer
%     timing information through queue send/receive latency. This is
%     discussed in Section VII.
